\section{Introducción}
\label{sec:introduccion}

Los modelos de lenguaje grandes responden generando texto: la secuencia de tokens \emph{es} la respuesta y
su longitud es el costo. En septiembre de 2026 apareció una familia distinta. Los \textbf{modelos de
decisión} (o modelos \emph{System One}) reciben una pregunta y un conjunto \textbf{cerrado y etiquetado} de
opciones, y responden con \textbf{una opción y una probabilidad}, sin emitir texto: terminan la inferencia
con un \emph{readout} de logits en la posición de cada etiqueta permitida \cite{jev2026blog,jev2026docs}.
Su costo de salida es nulo ---el proveedor no cobra los tokens de salida---, su latencia es de decenas de milisegundos, y
se ofrecen como el componente que decide \guillemotleft qué hacer\guillemotright{} dentro de un sistema
mayor: clasificar, enrutar, verificar, abstenerse.

Jev, de TypeSafe, abrió la categoría el 18 de septiembre de 2026 \cite{jev2026blog}. En menos de un mes
floreció un ecosistema \textbf{open-source} de pesos abiertos que replica su contrato: \textbf{Laya} (421M,
la única implementación que publica el bucle de entrenamiento RLCD completo) \cite{layaRepo}, \textbf{Kev}
como familia entrenable \cite{kevRepo}, y una decena de derivados (Nimble, SemIf, Jeff, jev-lite, AnyJev,
entre otros) \cite{nimble,semif,jeff,jevLite,anyjev}. Todos exponen el mismo campo:
\textbf{\texttt{confidence}}.

Ese campo es el que hace útil al modelo. Un despliegue típico no usa la etiqueta a secas: fija un
\textbf{umbral} sobre \texttt{confidence} y decide \emph{actuar}, o \emph{escalar a un humano} cuando el
modelo no está seguro \cite{chow1970,geifman2017}. El umbral, por tanto, \textbf{es} la política. Y aquí
está el problema: aunque el campo se llame igual en todo el ecosistema, \textbf{su definición no es la
misma}.

\begin{table}[t]
\centering
\caption{Definiciones de \texttt{confidence} presentes en el ecosistema.}
\label{tab:definiciones}
\footnotesize
\resizebox{\columnwidth}{!}{%
\begin{tabular}{@{}lll@{}}
\toprule
implementación & definición de \texttt{confidence} & fórmula \\
\midrule
Jev & concentración & $(n\,p_{\max} - 1)/(n - 1)$ \cite{jev2026docs} \\
Laya (\texttt{choice}, \texttt{score}) & entropía normalizada & $1 - H(p)/\log n$ \cite{layastudio} \\
Laya (\texttt{noul}) & probabilidad binaria & $\max(p, 1 - p)$ \cite{layastudio} \\
jev-lite, Laya \texttt{answer\_confidence} & probabilidad ganadora & $p_{\max}$ \cite{jevLite} \\
otros derivados & margen, prior aprendido & $p_{(1)} - p_{(2)}$, \dots \\
\bottomrule
\end{tabular}}
\end{table}

El efecto es inmediato: una sola distribución sobre tres opciones, $(0.0118, 0.4629, 0.5253)$, produce
\textbf{tres valores distintos} de \texttt{confidence}: 0.288 con la fórmula de Jev, 0.3199 con la entropía
de Laya y 0.5253 con $p_{\max}$ \cite{rohitghumare}. Es el mismo dato y el número cambia según quién lo
calcule. Un umbral calibrado sobre una definición, portado a otra, \textbf{no significa lo mismo}, y el
ecosistema los porta entre modelos como si fueran intercambiables.

La afirmación cualitativa \guillemotleft los umbrales no transfieren\guillemotright{} ya circula en
anécdotas \cite{jebadiah}. Lo que \textbf{no existe} es la medición sistemática: \textbf{cuánto} se degrada
la cobertura y la precisión del umbral al portarlo entre definiciones, modelos y datasets. Esa tabla es la
que este trabajo construye, y la construye con un instrumento que se regenera con un comando.

\textbf{Contribuciones.}

\begin{enumerate}
\item Una \textbf{especificación formal} de cada definición de \texttt{confidence} en uso, como función de
  la distribución devuelta y de la cardinalidad $n$ (dominio, rango, monotonía), con tres
  \textbf{proposiciones} que separan las definiciones que son \emph{equivalentes en orden} (transferencia
  exacta de umbral) de las que no lo son, y que muestran que la incompatibilidad aparece solo a partir de
  $n = 3$.
\item La librería \textbf{\texttt{jev-conf-metrics}}, que implementa todas las definiciones y las métricas de
  calibración y discriminación, con \textbf{67 tests de conformidad} contra salidas publicadas.
\item La \textbf{matriz de transferencia de umbrales} (definición origen $\times$ definición destino
  $\times$ modelo $\times$ dataset), el \textbf{dataset de 65\,185 mediciones crudas} con versión y hash por
  fila, las \textbf{figuras} y los \textbf{scripts regeneradores} que la reproducen sin GPU.
\item El \textbf{pre-registro} de la cobertura objetivo, las definiciones y los datasets antes de la corrida
  principal (DOI \texttt{10.5281/zenodo.23266795}).
\end{enumerate}

Con ese instrumento, el trabajo \textbf{mide} y responde: sobre 65\,185 decisiones medidas con cinco
configuraciones de modelo, la definición más alta y la más baja de la \textbf{misma} distribución difieren
con mediana \textbf{0.288}; y portar el umbral sin traducirlo se equivoca del orden de \textbf{cien veces}
más que portarlo con el mapa matemático cuando este existe.

El trabajo no reporta un hallazgo sobre si la calibración \guillemotleft viaja\guillemotright{} o no:
entrega el \textbf{instrumento con el que se mide}. Es el patrón de \cite{guo2017}, que no descubrió que las
redes modernas están des-calibradas, sino que dio la métrica y la corrección que el campo entero cita. La
aportación sobrevive al hueco: cualquiera que afirme que los umbrales no transfieren tendrá que
\textbf{citar la matriz}, no reemplazarla.

\textbf{Organización.} La Sección 2 sitúa el trabajo frente a la calibración, las reglas de puntuación
propias, la predicción selectiva y el ecosistema Jev-OSS. La Sección 3 fija el marco teórico: definiciones,
calibración frente a discriminación, las proposiciones de equivalencia en orden y la definición formal de la
transferencia de umbrales. Las Secciones 4 a 7 describen el método, los resultados, la discusión y las
conclusiones, y listan los artefactos entregados.
